% =====================================================================
%  Лекция 1.
%  Дата: 2026-09-24
% =====================================================================
\section{Основные понятия}
\label{lec:01}

\subsection{Выборочное пространство}

Рассмотрим случайную величину
\[
  \xi \colon \Omega \to \R.
\]
Рассмотрим вероятностное пространство
\[
  \bigl(\R,\, \mathcal{B}(\R),\, \P_{\xi}\bigr),
\]
где
\begin{itemize}
  \item $\mathcal{B}(\R)$~--- борелевская $\sigma$-алгебра на~$\R$;
  \item $\P_\xi(A)$~--- вероятностная мера такая, что
        $\P_\xi\bigl((-\infty, x]\bigr) = F(x) = \P(\xi \le x)$.
\end{itemize}

\medskip

Если говорить о наборе случайных величин $(\xi_1, \ldots, \xi_n)$, то вероятностное
пространство определяется как
\[
  \bigl(\R^n,\, \mathcal{B}(\R^n),\, \P_{\xi}\bigr),
\]
где
\begin{itemize}
  \item $\mathcal{B}(\R^n)$~--- борелевская $\sigma$-алгебра на~$\R^n$;
  \item $\P_\xi$~--- совместное распределение случайных величин $(\xi_1, \ldots, \xi_n)$:
        \[
          \P_\xi\bigl((-\infty, x_1] \times \cdots \times (-\infty, x_n]\bigr)
          = F(x_1, \ldots, x_n)
          = \P(\xi_1 \le x_1, \ldots, \xi_n \le x_n).
        \]
\end{itemize}

% \begin{definition}
%   \emph{Выборочным пространством} называется вероятностное пространство
%   $\bigl(\R^n,\, \mathcal{B}(\R^n),\, \P_{\xi}\bigr)$.
% \end{definition}

\begin{definition}
  \emph{Выборкой} $X_{[n]} = (X_1, \ldots, X_n)$ размера~$n$ называется совокупность
  взаимно независимых реализаций случайной величины~$\xi$, где $X_i$~--- числовая
  реализация случайной величины~$\xi_i$ в $i$-м эксперименте.
\end{definition}

Совместной функцией распределения
$F_{X_{[n]}}(x_1, \ldots, x_n) = F_\xi(x_1) \cdots F_\xi(x_n)$
однозначно определяется совместная вероятность $\P_{X_{[n]}}$.

\begin{definition}
  \emph{Генеральной совокупностью} называется множество всех возможных значений случайной величины $\xi$.
\end{definition}

Получили выборочное пространство:
\begin{definition}
  \emph{Выборочным пространством} называется вероятностное пространство
  $\bigl(\R^n,\, \mathcal{B}(\R^n),\, \P_{X_{[n]}}\bigr)$.
\end{definition}

\begin{reminder}
  Функция $f \colon \R^n \to \R^m$ называется \emph{борелевской}, если
  $\forall B \in \mathcal{B}(\R^m)$ выполнено $f^{-1}(B) \in \mathcal{B}(\R^n)$.
\end{reminder}

\begin{definition}
  \emph{Статистикой} называется любая борелевская функция, заданная на
  выборочном пространстве.
\end{definition}

\subsection{Эмпирическая функция распределения. Гистограмма}

Пусть задана генеральная совокупность случайной величины~$\xi$
и выборка размера~$n$: $X_{[n]} = (X_1, \ldots, X_n)$.
\begin{definition}
  \emph{Эмпирическим распределением} называется вероятностная мера,
  определённая следующим образом:
  \[
    \Pemp(B) = \P_{X_{[n]}}^{*}(B) = \frac{\nu(B)}{n},
  \]
  где $\nu(B)$~--- количество элементов выборки, попавших в борелевское множество
  $B \in \mathcal{B}(\R)$.
\end{definition}

\begin{definition}
  \emph{Эмпирической функцией распределения} называется функция
  \[
    F_n^{*}(x) = \Pemp\bigl((-\infty, x]\bigr)
    = \frac{\nu\bigl((-\infty, x]\bigr)}{n}, \qquad x \in \R.
  \]
\end{definition}

Введём порядковые статистики:
\begin{align*}
  X_{(1)} &= \min(X_1, \ldots, X_n), \\
  X_{(2)} &= \min\bigl(\{X_1, \ldots, X_n\} \setminus \{X_{(1)}\}\bigr), \\
          &\;\;\vdots \\
  X_{(n)} &= \max(X_1, \ldots, X_n).
\end{align*}
Очевидно, что $X_{(1)} \le X_{(2)} \le \ldots \le X_{(n)}$.

\begin{definition}
  \emph{Вариационным рядом} называется последовательность порядковых статистик
  \[
    X_{(1)} \le X_{(2)} \le \ldots \le X_{(n)}.
  \]
\end{definition}

Если предположить, что все элементы вариационного ряда различны,
то есть $X_{(1)} < X_{(2)} < \ldots < X_{(n)}$, то эмпирическую
функцию распределения можно построить следующим образом:
\begin{equation}
  \label{eq:ecdf}
  F_n^*(x) =
  \begin{cases}
    0,           & \text{если } x < X_{(1)}; \\
    \frac{1}{n}, & \text{если } X_{(1)} \le x < X_{(2)}; \\
    \frac{2}{n}, & \text{если } X_{(2)} \le x < X_{(3)}; \\
    \ldots \\
    \frac{k}{n}, & \text{если } X_{(k)} \le x < X_{(k+1)}; \\
    \ldots \\
    1,           & \text{если } x \ge X_{(n)}.
  \end{cases}
\end{equation}

Теперь можно построить гистограмму. Возьмём интервал $(a, b)$, где $a < X_{(1)}$ и $b > X_{(n)}$,
и разобьём его точками
$a = a_0 < a_1 < \ldots < a_m = b$, где $m \le n$.
Пусть $\nu_i$~--- количество элементов выборки, попавших в интервал $(a_{i-1}, a_{i}]$.
Тогда $\nu_1 + \nu_2 + \ldots + \nu_m = n$. Положим
\[
  l_i = a_i - a_{i-1}, \qquad h_i = \frac{\nu_i}{n l_i}, \qquad i = 1, \ldots, m.
\]

Получаем гистограмму:
\begin{equation}
  \label{eq:histogram}
  f_n^{*}(x) =
  \begin{cases}
    h_1, & \text{если } x \in (a_0, a_1]; \\
    h_2, & \text{если } x \in (a_1, a_2]; \\
    \ldots \\
    h_m, & \text{если } x \in (a_{m-1}, a_m].
  \end{cases}
\end{equation}

Гистограмма $f_n^{*}(x)$~--- эмпирический аналог плотности распределения. Диаграмму
относительных частот можно построить аналогично, убрав при вычислении $h_i$ деление на $l_i$.
Если в гистограмме убрать $n$, то получим гистограмму частот.
